% !PRINTED 275
\weberpage{275}
Here $\gamma$ denotes a primitive root, hence $\gamma^{\frac{p^2-1}{\nu}}$
a primitive $\nu$th root of unity of the congruence field
$\mathfrak{E}_{2,p}$ (\S.\,67).

3. \gesperrt{Both poles of a substitution $\Theta$ are of equal multiplicity}
(\S.\,52, 4.).

If $Q$ is the cyclic group of degree $\nu$ of the powers
of $\Theta$, and $n = \nu\mu$, then one obtains
\[
G = Q + \psi_1 Q + \cdots + \psi_{\mu-1} Q,
\]
where $\psi_1, \ldots, \psi_{\mu-1}$ are substitutions from $G$, and

4. \gesperrt{the numbers}
\weq{4}{\alpha,\ \psi_1(\alpha) = \alpha_1,\ \psi_2(\alpha) = \alpha_2,\ \ldots,\ \psi_{\mu-1}(\alpha) = \alpha_{\mu-1}}
\gesperrt{form a system of conjugate poles of equal multiplicity of the group $G$.}

All the poles of the group $G$ can thus be gathered into systems
of conjugate poles, and in this way we obtain the indeterminate
equation
\weq{5}{\begin{aligned}
2n - 2 &= \mu(\nu-1) + \mu'(\nu'-1) + \mu''(\nu''-1) \ldots \\
&= nh - \mu - \mu' - \mu'' \ldots,
\end{aligned}}
where $h$ is the number of systems of conjugate poles; and of this
equation we have seen in \S.\,52 that it admits only a limited
number of solutions.

In order to find the divisors of the congruence groups corresponding
to these solutions, we may simply put, in the formulae of the
polyhedral groups set up in the eighth chapter, for the roots of
unity occurring there the roots of unity of the field
$\mathfrak{E}_{2,p}$ defined in \S.\,67, with which indeed
calculation proceeds by the same rules. In this, of course, only such
roots of unity can occur whose degree is a divisor of $p^2 - 1$.

Finally, for every such group it remains to be investigated whether
its substitutions are contained in $L_p$ or in a group isomorphic
with it. We then have the following cases of solutions of
equation (5), where $\gamma$ always denotes a primitive root of the field
$\mathfrak{E}_{2,p}$ (\S.\,55).

\gesperrt{I. Cyclic groups of degree $n$,}\quad
$\begin{aligned} h &= 2,\ \nu = \nu' = n \\ \mu &= \mu' = 1 \end{aligned}$
\[
C_n = \begin{pmatrix} \gamma^{\frac{p^2-1}{2n}\lambda}, & 0 \\[2pt] 0, & \gamma^{-\frac{(p^2-1)}{2n}\lambda} \end{pmatrix},
\quad \lambda = 0, 1, \ldots, n-1.
\]
